# 1. Introduction

Nuclear fission is intrinsically a quantum-mechanical process. The formation of an excited compound nucleus, the evolution of nuclear degrees of freedom toward scission, quantum tunnelling through the fission barrier, and the subsequent formation of fragments are all governed, at the microscopic level, by quantum dynamics. Modern microscopic approaches to fission therefore describe the nucleus as a complex interacting many-body quantum system and increasingly employ nuclear density functional theory, time-dependent density functional methods, and collective quantum dynamics to investigate the path from excitation to scission.

At the reactor-engineering scale, however, this microscopic quantum dynamics is not treated as an active degree of freedom. Reactor physics is formulated primarily in terms of neutron transport, reaction rates, isotopic concentrations, thermal-hydraulic variables and macroscopic feedback mechanisms. The microscopic nuclear dynamics is effectively incorporated into quantities such as reaction cross sections, decay constants, neutron multiplicities and energy-dependent nuclear data. Consequently, the internal quantum state of the fissile nuclei is normally not regarded as a controllable variable of the reactor.

Reactor control therefore acts predominantly on the environment in which fission takes place rather than on the quantum state from which fission proceeds. Control rods modify neutron absorption; moderators alter the neutron energy distribution; soluble absorbers modify macroscopic absorption; coolant and fuel temperatures provide reactivity feedback; and changes in geometry or composition affect neutron leakage and multiplication. In each case, the fundamental strategy is to control the neutron population or its interaction with an effectively passive fissile medium.

This situation differs markedly from that encountered in atomic and optical physics. The development of coherent control has made it possible not only to predict transition probabilities but also to manipulate quantum states by controlling the amplitude, frequency, timing and phase of an applied electromagnetic field. The laser is perhaps the most familiar technological manifestation of this principle: rather than treating atomic transition probabilities as immutable statistical parameters, the state of the active medium and its interaction with the radiation field become part of the dynamics of the device.

Historically, extending such concepts to nuclear degrees of freedom appeared extremely challenging because typical nuclear transition energies are many orders of magnitude larger than those encountered in conventional optical systems. Nevertheless, the field of nuclear quantum optics has progressively demonstrated that the interaction between coherent electromagnetic radiation and nuclear states can itself become experimentally accessible. Early theoretical studies showed that sufficiently intense high-frequency radiation could enable resonant laser–nucleus interactions and explicitly proposed the preparation, control and detection of nuclear states as a long-term objective of nuclear quantum optics.

Subsequent developments have moved this concept beyond purely theoretical considerations. Coherent control of nuclear excitations has been experimentally demonstrated using X-ray fields. In particular, coherent manipulation of collective nuclear excitations in (^{57}\mathrm{Fe}) has shown that interference and phase-sensitive techniques familiar from quantum optics can be transferred to nuclear resonances. More recently, direct laser spectroscopy of the exceptionally low-energy nuclear isomeric transition in (^{229}\mathrm{Th}) has provided another striking demonstration that a nuclear transition can be directly addressed using optical techniques.

At substantially higher excitation energies, theoretical investigations of intense coherent MeV photon fields have also considered nuclear excitation in regimes in which photon absorption competes with intrinsic nuclear equilibration. Such models explicitly include dipole absorption, stimulated emission, neutron evaporation and induced fission, showing that coherent radiation fields can, at least theoretically, participate in the dynamical evolution of highly excited heavy nuclei.

At the same time, microscopic theories of induced fission are increasingly capable of following the time evolution of the fissioning system itself. Rather than representing fission exclusively through an averaged cross section, these approaches study the evolution of nuclear configurations and collective wave packets toward scission, including fluctuations, dissipation, pairing dynamics and the distribution of possible fission outcomes.

These two developments suggest a question that is distinct from both conventional reactor physics and conventional nuclear quantum optics.

**Can externally controlled preparation or manipulation of nuclear states in fissile matter modify the subsequent probability or dynamics of fission in a way that can be represented as a controllable physical variable?**

The distinction is important. The question is not whether electromagnetic radiation can induce fission; photofission is well established. Nor is it simply whether fission must ultimately be described quantum mechanically. Rather, the question considered here is whether the *state dependence* of the fission process can itself be exploited as a control mechanism.

Conventionally, a fission rate may schematically be represented as

[
R_f = N_f \int \sigma_f(E),\Phi(E),dE ,
]

where (N_f) is the number of fissile nuclei, (\Phi(E)) is the neutron flux and (\sigma_f(E)) is the corresponding fission cross section. Within this description, reactor control primarily modifies (\Phi(E)), while (\sigma_f(E)) is treated as a property supplied by nuclear data.

The conceptual extension investigated in this work is to consider whether the effective fission response could instead depend explicitly on an externally prepared nuclear state,

[
\sigma_f
\rightarrow
\sigma_f(E,\rho_N),
]

where (\rho_N) denotes the nuclear quantum state, or more generally the relevant nuclear density operator or state distribution. If (\rho_N) can be influenced by an external coherent interaction,

[
u_Q(t)
\longrightarrow
\rho_N(t)
\longrightarrow
\sigma_f^{\mathrm{eff}}
\longrightarrow
R_f ,
]

then the microscopic state of the fissile material becomes, in principle, an additional degree of freedom in the control of a fission system.

This possibility should not be interpreted as the proposal of an immediately realizable “quantum nuclear reactor”. Significant physical obstacles are apparent, including the extremely rapid intrinsic equilibration of highly excited nuclei, decoherence, the energy scales required to address relevant nuclear excitations, the complexity of the fission continuum, and the challenge of producing a sufficiently large controlled effect in macroscopic matter. Existing studies of strong coherent photon–nucleus interactions already indicate the importance of competition between externally driven excitation and rapid internal nuclear equilibration.

The purpose of the present work is therefore more limited and more fundamental. We ask whether a consistent theoretical framework can be constructed in which nuclear-state preparation constitutes an independent control variable for fission, and under which physical conditions such control could produce a non-negligible modification of fission observables.

In this sense, the central hypothesis can be stated succinctly:

[
\boxed{
\text{The quantum state of fissile matter, normally implicit in nuclear reaction data, may in principle become an explicit control variable.}
}
]

Establishing whether this hypothesis is physically meaningful requires connecting three areas that are normally treated separately: microscopic quantum theories of nuclear fission, coherent control of nuclear excitations, and macroscopic reactor dynamics. The objective of this work is to construct the first conceptual bridge between these descriptions and to identify the minimal physical conditions required for quantum-state control to influence a macroscopic fission system.
